\documentclass[11pt,a4paper]{article}
\usepackage[utf8]{inputenc}
\usepackage[T1]{fontenc}
\usepackage[spanish,es-minimal,es-noshorthands]{babel}
\AtBeginDocument{\renewcommand{\sin}{\operatorname{sen}}\renewcommand{\arcsin}{\operatorname{arcsen}}}
\usepackage{lmodern}
\usepackage[sfdefault,lining]{FiraSans}
\usepackage[a4paper,top=2.3cm,bottom=2.2cm,left=1.8cm,right=1.8cm,headheight=1.3cm,headsep=0.4cm,footskip=1cm]{geometry}
\usepackage{amsmath,amssymb}
\usepackage{icomma}
\usepackage{graphicx}
\usepackage[table]{xcolor}
\usepackage{tikz}
\usepackage{fontawesome5}
\usepackage{ccicons}
\usepackage{fancyhdr}
\usepackage{lastpage}
\usepackage{tasks}
\usepackage{needspace}
\usepackage[most]{tcolorbox}
\usepackage{array}
\usepackage[hidelinks]{hyperref}
\AtBeginDocument{\hypersetup{pdftitle={Ficha de ejercicios de la unidad 2: Números reales. 4º ESO},pdfsubject={Matemáticas A 4º ESO},pdfkeywords={números reales, números racionales, números irracionales, fracción generatriz, intervalos, semirrectas, unión, intersección, raíces, radicales, potencias de exponente fraccionario, radicales equivalentes, extracción de factores, operaciones con radicales, racionalización}}}

\definecolor{navy}{HTML}{1B2748}
\definecolor{accent}{HTML}{E8603A}
\definecolor{teal}{HTML}{1F9E90}
\definecolor{soft}{HTML}{EEF1F8}
\definecolor{muted}{HTML}{5D6685}

\setlength{\parindent}{0pt}
\setlength{\parskip}{3pt}
\settasks{label=\alph*), label-format=\bfseries\color{navy}, label-width=1.3em, item-indent=1.8em, column-sep=0.6em, after-item-skip=5pt, before-skip=3pt, after-skip=4pt}

% ---------- capçalera i peu ----------
\newcommand{\llicencia}{\href{https://creativecommons.org/licenses/by-nc-sa/4.0/deed.es}{\ccbyncsa}\ \ Licencia Creative Commons Atribución-NoComercial-CompartirIgual 4.0 Internacional (CC BY-NC-SA 4.0)}
\pagestyle{fancy}
\fancyhf{}
\fancyhead[L]{\raisebox{-0.35\height}{\includegraphics[height=0.95cm]{logo-gva.png}}}
\fancyhead[C]{\small\color{muted}Ficha de ejercicios · Unidad 2: Números reales · 4\textsuperscript{o} ESO A}
\fancyhead[R]{\raisebox{-0.35\height}{\includegraphics[height=0.95cm]{logo-vermellar.png}}}
\fancyfoot[L]{\scriptsize\color{muted}\llicencia}
\fancyfoot[R]{\small\color{muted}\thepage}
\renewcommand{\headrulewidth}{0.4pt}
\renewcommand{\headrule}{\color{muted!40}\hrule height \headrulewidth}
\fancypagestyle{primera}{\fancyhf{}\renewcommand{\headrulewidth}{0pt}
  \fancyfoot[L]{\scriptsize\color{muted}\llicencia}
  \fancyfoot[R]{\small\color{muted}\thepage}}

% ---------- nivells i exercicis ----------
\newcommand{\estrelles}[1]{{\color{accent}\ifcase#1\or\faStar\faStar[regular]\faStar[regular]\or\faStar\faStar\faStar[regular]\or\faStar\faStar\faStar\fi}}
\newcounter{exer}
\newcommand{\ex}[1]{\par\needspace{4\baselineskip}\medskip\refstepcounter{exer}\noindent{\color{navy}\bfseries\theexer.}\ \ {\small\estrelles{#1}}\ \ }

\newcommand{\seccio}[2]{\par\needspace{16\baselineskip}\bigskip
  \begin{tcolorbox}[colback=navy,colframe=navy,coltext=white,arc=2mm,boxrule=0pt,left=3mm,right=3mm,top=1.2mm,bottom=1.2mm]
  {\large\bfseries #1\quad #2}
  \end{tcolorbox}\vspace{-2pt}}

\newcommand{\exf}[1]{\par\medskip\refstepcounter{exer}\noindent{\color{navy}\bfseries\theexer.}\ \ {\small\estrelles{#1}}\ \ }
\newcolumntype{C}[1]{>{\centering\arraybackslash}p{#1}}
% quadre de recordatori (només a la fitxa d'exercicis)
\newcommand{\recorda}[1]{\par\smallskip\begin{tcolorbox}[colback=accent!7,colframe=accent!55,arc=2mm,boxrule=0.5pt,left=2.5mm,right=2.5mm,top=1mm,bottom=1mm]\small{\color{accent}\faLightbulb[regular]}\ \textbf{Recuerda:}\ #1\end{tcolorbox}}
\newcommand{\res}[1]{#1}\newcommand{\rest}[1]{#1}

\def\fmneed{11}% que el títol d'una secció no quede sol al final del full
% =====================================================================
% Capçalera FentMates de les fitxes: un color per a cada curs.
% Els colors són els de la pàgina inicial (LEVELS a assets/js/materials.js);
% el logo ve de _eines/logo/pdf/ (python3 _eines/logo/genera.py).
%
% Ús, al final del preàmbul de la unitat:
%   % =====================================================================
% Capçalera FentMates de les fitxes: un color per a cada curs.
% Els colors són els de la pàgina inicial (LEVELS a assets/js/materials.js);
% el logo ve de _eines/logo/pdf/ (python3 _eines/logo/genera.py).
%
% Ús, al final del preàmbul de la unitat:
%   % =====================================================================
% Capçalera FentMates de les fitxes: un color per a cada curs.
% Els colors són els de la pàgina inicial (LEVELS a assets/js/materials.js);
% el logo ve de _eines/logo/pdf/ (python3 _eines/logo/genera.py).
%
% Ús, al final del preàmbul de la unitat:
%   \input{../comu/fentmates.tex}
%   \fmcurs{1eso}
% i, al començament del document:
%   \capcalera{Matemàtiques}{1r ESO}{Fitxa d'exercicis · Unitat 2}{Divisibilitat}
% Opcional abans de \input: \def\fmneed{14} (espai mínim abans de cada secció).
% =====================================================================
\makeatletter
\def\fm@logodir{../../logo/pdf}
\providecommand{\fmneed}{10}
\def\fm@cursnom{}

% clau del curs, color i color del text damunt del color
\newcommand{\fm@curs}[3]{\@namedef{fm@color@#1}{#2}\@namedef{fm@text@#1}{#3}}
\fm@curs{1eso}{E6B84A}{navy}
\fm@curs{2eso}{3B82F6}{white}
\fm@curs{3eso}{FF7A50}{white}
\fm@curs{4eso-a}{A78BFA}{white}
\fm@curs{4eso-b}{6D28D9}{white}
\fm@curs{1batx-mi}{9B1C1C}{white}
\fm@curs{1batx-ccss}{EF6B6B}{white}
\fm@curs{2batx-mii}{1E6B47}{white}
\fm@curs{2batx-ccss}{2BB3A3}{white}

\newcommand{\fmcurs}[1]{%
  \@ifundefined{fm@color@#1}{\PackageError{fentmates}{Curs desconegut: #1}{}}{}%
  \edef\fm@tmp{\noexpand\definecolor{curs}{HTML}{\@nameuse{fm@color@#1}}%
               \noexpand\colorlet{curstext}{\@nameuse{fm@text@#1}}}\fm@tmp
  \colorlet{cursfosc}{curs!72!black}%
  \def\fm@logo{\fm@logodir/vertical-#1.pdf}%
  \def\fm@marca{\fm@logodir/marca-#1.pdf}}

% ---------- capçalera de la primera pàgina ----------
% #1 matèria (petit), #2 curs, #3 tipus de document i unitat, #4 títol
\newcommand{\capcalera}[4]{%
  \gdef\fm@cursnom{#2}%
  \par\noindent
  \begin{tcolorbox}[colback=curs,colframe=curs,boxrule=0pt,arc=2.2mm,width=3.5cm,height=2.15cm,
      valign=center,left=3mm,right=2mm,top=1mm,bottom=1mm,nobeforeafter,box align=center]
    \color{curstext}{\footnotesize #1}\par\vspace{1pt}{\fontsize{19}{21}\selectfont\bfseries #2\par}%
  \end{tcolorbox}\hspace{5mm}%
  \begin{minipage}[c]{\dimexpr\linewidth-3.5cm-5mm-2.9cm\relax}
    {\color{cursfosc}\small\bfseries\MakeUppercase{#3}}\par\vspace{2pt}
    {\color{navy}\fontsize{20}{23}\selectfont\bfseries #4\par}
  \end{minipage}\hfill
  \begin{minipage}[c]{2.6cm}\raggedleft\includegraphics[width=2.6cm]{\fm@logo}\end{minipage}\par
  \vspace{2pt}}

% ---------- barres de secció del color del curs ----------
\renewcommand{\seccio}[2]{\par\needspace{\fmneed\baselineskip}\bigskip
  \begin{tcolorbox}[colback=curs,colframe=curs,coltext=curstext,arc=2mm,boxrule=0pt,left=3mm,right=3mm,top=1.2mm,bottom=1.2mm]
  {\large\bfseries #1\quad #2}
  \end{tcolorbox}\vspace{-2pt}}

% ---------- estrelles de dificultat del color del curs (a Batxillerat no se'n posen) ----------
\@ifundefined{estrelles}{}{%
  \renewcommand{\estrelles}[1]{{\color{curs}\ifcase#1\or\faStar\faStar[regular]\faStar[regular]\or\faStar\faStar\faStar[regular]\or\faStar\faStar\faStar\fi}}}

% ---------- capçalera i peu de les pàgines ----------
% primera pàgina: logos institucionals al peu
\fancypagestyle{primera}{\fancyhf{}\renewcommand{\headrulewidth}{0pt}
  \fancyfoot[L]{\raisebox{-0.32\height}{\includegraphics[height=0.78cm]{logo-gva.png}}\hspace{5mm}%
    \raisebox{-0.32\height}{\includegraphics[height=0.78cm]{logo-vermellar.png}}}
  \fancyfoot[R]{\scriptsize\color{muted}\href{https://creativecommons.org/licenses/by-nc-sa/4.0/deed.ca}{\ccbyncsa}\ \ CC BY-NC-SA 4.0 · fentmates.org\qquad\small\thepage}}
% la resta: el símbol a l'esquerra i el curs a la dreta, amb el títol al centre com abans
\fancyhead[L]{\raisebox{-0.3\height}{\includegraphics[height=0.62cm]{\fm@marca}}}
\fancyhead[R]{\small\bfseries\color{cursfosc}\fm@cursnom}
\renewcommand{\headrule}{\color{curs!45}\hrule height \headrulewidth}
\makeatother

%   \fmcurs{1eso}
% i, al començament del document:
%   \capcalera{Matemàtiques}{1r ESO}{Fitxa d'exercicis · Unitat 2}{Divisibilitat}
% Opcional abans de \input: \def\fmneed{14} (espai mínim abans de cada secció).
% =====================================================================
\makeatletter
\def\fm@logodir{../../logo/pdf}
\providecommand{\fmneed}{10}
\def\fm@cursnom{}

% clau del curs, color i color del text damunt del color
\newcommand{\fm@curs}[3]{\@namedef{fm@color@#1}{#2}\@namedef{fm@text@#1}{#3}}
\fm@curs{1eso}{E6B84A}{navy}
\fm@curs{2eso}{3B82F6}{white}
\fm@curs{3eso}{FF7A50}{white}
\fm@curs{4eso-a}{A78BFA}{white}
\fm@curs{4eso-b}{6D28D9}{white}
\fm@curs{1batx-mi}{9B1C1C}{white}
\fm@curs{1batx-ccss}{EF6B6B}{white}
\fm@curs{2batx-mii}{1E6B47}{white}
\fm@curs{2batx-ccss}{2BB3A3}{white}

\newcommand{\fmcurs}[1]{%
  \@ifundefined{fm@color@#1}{\PackageError{fentmates}{Curs desconegut: #1}{}}{}%
  \edef\fm@tmp{\noexpand\definecolor{curs}{HTML}{\@nameuse{fm@color@#1}}%
               \noexpand\colorlet{curstext}{\@nameuse{fm@text@#1}}}\fm@tmp
  \colorlet{cursfosc}{curs!72!black}%
  \def\fm@logo{\fm@logodir/vertical-#1.pdf}%
  \def\fm@marca{\fm@logodir/marca-#1.pdf}}

% ---------- capçalera de la primera pàgina ----------
% #1 matèria (petit), #2 curs, #3 tipus de document i unitat, #4 títol
\newcommand{\capcalera}[4]{%
  \gdef\fm@cursnom{#2}%
  \par\noindent
  \begin{tcolorbox}[colback=curs,colframe=curs,boxrule=0pt,arc=2.2mm,width=3.5cm,height=2.15cm,
      valign=center,left=3mm,right=2mm,top=1mm,bottom=1mm,nobeforeafter,box align=center]
    \color{curstext}{\footnotesize #1}\par\vspace{1pt}{\fontsize{19}{21}\selectfont\bfseries #2\par}%
  \end{tcolorbox}\hspace{5mm}%
  \begin{minipage}[c]{\dimexpr\linewidth-3.5cm-5mm-2.9cm\relax}
    {\color{cursfosc}\small\bfseries\MakeUppercase{#3}}\par\vspace{2pt}
    {\color{navy}\fontsize{20}{23}\selectfont\bfseries #4\par}
  \end{minipage}\hfill
  \begin{minipage}[c]{2.6cm}\raggedleft\includegraphics[width=2.6cm]{\fm@logo}\end{minipage}\par
  \vspace{2pt}}

% ---------- barres de secció del color del curs ----------
\renewcommand{\seccio}[2]{\par\needspace{\fmneed\baselineskip}\bigskip
  \begin{tcolorbox}[colback=curs,colframe=curs,coltext=curstext,arc=2mm,boxrule=0pt,left=3mm,right=3mm,top=1.2mm,bottom=1.2mm]
  {\large\bfseries #1\quad #2}
  \end{tcolorbox}\vspace{-2pt}}

% ---------- estrelles de dificultat del color del curs (a Batxillerat no se'n posen) ----------
\@ifundefined{estrelles}{}{%
  \renewcommand{\estrelles}[1]{{\color{curs}\ifcase#1\or\faStar\faStar[regular]\faStar[regular]\or\faStar\faStar\faStar[regular]\or\faStar\faStar\faStar\fi}}}

% ---------- capçalera i peu de les pàgines ----------
% primera pàgina: logos institucionals al peu
\fancypagestyle{primera}{\fancyhf{}\renewcommand{\headrulewidth}{0pt}
  \fancyfoot[L]{\raisebox{-0.32\height}{\includegraphics[height=0.78cm]{logo-gva.png}}\hspace{5mm}%
    \raisebox{-0.32\height}{\includegraphics[height=0.78cm]{logo-vermellar.png}}}
  \fancyfoot[R]{\scriptsize\color{muted}\href{https://creativecommons.org/licenses/by-nc-sa/4.0/deed.ca}{\ccbyncsa}\ \ CC BY-NC-SA 4.0 · fentmates.org\qquad\small\thepage}}
% la resta: el símbol a l'esquerra i el curs a la dreta, amb el títol al centre com abans
\fancyhead[L]{\raisebox{-0.3\height}{\includegraphics[height=0.62cm]{\fm@marca}}}
\fancyhead[R]{\small\bfseries\color{cursfosc}\fm@cursnom}
\renewcommand{\headrule}{\color{curs!45}\hrule height \headrulewidth}
\makeatother

%   \fmcurs{1eso}
% i, al començament del document:
%   \capcalera{Matemàtiques}{1r ESO}{Fitxa d'exercicis · Unitat 2}{Divisibilitat}
% Opcional abans de \input: \def\fmneed{14} (espai mínim abans de cada secció).
% =====================================================================
\makeatletter
\def\fm@logodir{../../logo/pdf}
\providecommand{\fmneed}{10}
\def\fm@cursnom{}

% clau del curs, color i color del text damunt del color
\newcommand{\fm@curs}[3]{\@namedef{fm@color@#1}{#2}\@namedef{fm@text@#1}{#3}}
\fm@curs{1eso}{E6B84A}{navy}
\fm@curs{2eso}{3B82F6}{white}
\fm@curs{3eso}{FF7A50}{white}
\fm@curs{4eso-a}{A78BFA}{white}
\fm@curs{4eso-b}{6D28D9}{white}
\fm@curs{1batx-mi}{9B1C1C}{white}
\fm@curs{1batx-ccss}{EF6B6B}{white}
\fm@curs{2batx-mii}{1E6B47}{white}
\fm@curs{2batx-ccss}{2BB3A3}{white}

\newcommand{\fmcurs}[1]{%
  \@ifundefined{fm@color@#1}{\PackageError{fentmates}{Curs desconegut: #1}{}}{}%
  \edef\fm@tmp{\noexpand\definecolor{curs}{HTML}{\@nameuse{fm@color@#1}}%
               \noexpand\colorlet{curstext}{\@nameuse{fm@text@#1}}}\fm@tmp
  \colorlet{cursfosc}{curs!72!black}%
  \def\fm@logo{\fm@logodir/vertical-#1.pdf}%
  \def\fm@marca{\fm@logodir/marca-#1.pdf}}

% ---------- capçalera de la primera pàgina ----------
% #1 matèria (petit), #2 curs, #3 tipus de document i unitat, #4 títol
\newcommand{\capcalera}[4]{%
  \gdef\fm@cursnom{#2}%
  \par\noindent
  \begin{tcolorbox}[colback=curs,colframe=curs,boxrule=0pt,arc=2.2mm,width=3.5cm,height=2.15cm,
      valign=center,left=3mm,right=2mm,top=1mm,bottom=1mm,nobeforeafter,box align=center]
    \color{curstext}{\footnotesize #1}\par\vspace{1pt}{\fontsize{19}{21}\selectfont\bfseries #2\par}%
  \end{tcolorbox}\hspace{5mm}%
  \begin{minipage}[c]{\dimexpr\linewidth-3.5cm-5mm-2.9cm\relax}
    {\color{cursfosc}\small\bfseries\MakeUppercase{#3}}\par\vspace{2pt}
    {\color{navy}\fontsize{20}{23}\selectfont\bfseries #4\par}
  \end{minipage}\hfill
  \begin{minipage}[c]{2.6cm}\raggedleft\includegraphics[width=2.6cm]{\fm@logo}\end{minipage}\par
  \vspace{2pt}}

% ---------- barres de secció del color del curs ----------
\renewcommand{\seccio}[2]{\par\needspace{\fmneed\baselineskip}\bigskip
  \begin{tcolorbox}[colback=curs,colframe=curs,coltext=curstext,arc=2mm,boxrule=0pt,left=3mm,right=3mm,top=1.2mm,bottom=1.2mm]
  {\large\bfseries #1\quad #2}
  \end{tcolorbox}\vspace{-2pt}}

% ---------- estrelles de dificultat del color del curs (a Batxillerat no se'n posen) ----------
\@ifundefined{estrelles}{}{%
  \renewcommand{\estrelles}[1]{{\color{curs}\ifcase#1\or\faStar\faStar[regular]\faStar[regular]\or\faStar\faStar\faStar[regular]\or\faStar\faStar\faStar\fi}}}

% ---------- capçalera i peu de les pàgines ----------
% primera pàgina: logos institucionals al peu
\fancypagestyle{primera}{\fancyhf{}\renewcommand{\headrulewidth}{0pt}
  \fancyfoot[L]{\raisebox{-0.32\height}{\includegraphics[height=0.78cm]{logo-gva.png}}\hspace{5mm}%
    \raisebox{-0.32\height}{\includegraphics[height=0.78cm]{logo-vermellar.png}}}
  \fancyfoot[R]{\scriptsize\color{muted}\href{https://creativecommons.org/licenses/by-nc-sa/4.0/deed.ca}{\ccbyncsa}\ \ CC BY-NC-SA 4.0 · fentmates.org\qquad\small\thepage}}
% la resta: el símbol a l'esquerra i el curs a la dreta, amb el títol al centre com abans
\fancyhead[L]{\raisebox{-0.3\height}{\includegraphics[height=0.62cm]{\fm@marca}}}
\fancyhead[R]{\small\bfseries\color{cursfosc}\fm@cursnom}
\renewcommand{\headrule}{\color{curs!45}\hrule height \headrulewidth}
\makeatother

\fmcurs{4eso-a}

\begin{document}
\thispagestyle{primera}

\capcalera{Matemáticas A}{4\textsuperscript{o} ESO}{Ficha de ejercicios · Unidad 2}{\raggedright Números reales}
\vspace{10pt}
Nombre y apellidos: \rule{6.6cm}{0.4pt}\hfill Grupo: \rule{1.6cm}{0.4pt}\hfill Fecha: \rule{2.2cm}{0.4pt}
\vspace{8pt}

\begin{tcolorbox}[colback=soft,colframe=soft,arc=2mm,left=3mm,right=3mm,top=1.5mm,bottom=1.5mm]
\small Ejercicios graduados por niveles de dificultad:\quad \estrelles{1}\ básico \qquad \estrelles{2}\ medio \qquad \estrelles{3}\ avanzado.\\[2pt] Resuélvelos en el cuaderno mostrando todos los pasos. Usa la calculadora solo cuando el enunciado lo diga.
\end{tcolorbox}

% =====================================================
\seccio{1}{Clasificación de los números reales}

\exf{1} Indica el menor conjunto numérico ($\mathbb N$, $\mathbb Z$, $\mathbb Q$ o $\mathbb I$) al que pertenece cada número:\begin{tasks}(4)
  \task $\dfrac{18}{3}$
  \task $-\sqrt{25}$
  \task $0,\widehat7$
  \task $\sqrt{20}$
  \task $-\dfrac74$
  \task $\sqrt[3]{-27}$
  \task $\pi-1$
  \task $\sqrt{-9}$
  \task $2,\widehat{03}$
  \task $\sqrt{\dfrac49}$
  \task $1,0100100010\ldots$
  \task $\dfrac{\sqrt{12}}{\sqrt3}$
\end{tasks}\recorda{Primero haz las operaciones que se puedan hacer: $\frac{18}{3}=6$, $\sqrt[3]{-27}=-3$\ldots{} Un decimal exacto o periódico es racional; un decimal infinito no periódico es irracional.}

\ex{1} Copia la tabla y marca los conjuntos a los que pertenece cada número:\begin{center}\renewcommand{\arraystretch}{1.6}\begin{tabular}{|C{1.6cm}|C{1.1cm}|C{1.1cm}|C{1.1cm}|C{1.1cm}|C{1.1cm}|}\hline
\rowcolor{soft} & $\mathbb N$ & $\mathbb Z$ & $\mathbb Q$ & $\mathbb I$ & $\mathbb R$\\\hline
$-3$ &  &  &  &  & \\\hline
$\dfrac52$ &  &  &  &  & \\\hline
$\sqrt7$ &  &  &  &  & \\\hline
$0$ &  &  &  &  & \\\hline
$-1,\widehat6$ &  &  &  &  & \\\hline
$\sqrt{16}$ &  &  &  &  & \\\hline
$\dfrac{\pi}{3}$ &  &  &  &  & \\\hline
\end{tabular}\end{center}

\ex{1} Di si son verdaderas o falsas las afirmaciones siguientes. Si son falsas, pon un contraejemplo:\begin{tasks}(1)
  \task Todo número natural es entero.
  \task Todo número entero es natural.
  \task Todo decimal periódico es racional.
  \task $\sqrt2=1,414213562$, porque es lo que dice la calculadora.
  \task Hay números racionales que también son irracionales.
  \task La suma de dos irracionales siempre es irracional.
\end{tasks}

\ex{2} Escribe:\begin{tasks}(2)
  \task un número entero que no sea natural;
  \task un número racional que no sea entero;
  \task tres números irracionales;
  \task un racional y un irracional comprendidos entre 3 y 4.
\end{tasks}

\ex{2} Ordena de menor a mayor e indica cuáles son racionales y cuáles irracionales (puedes usar la calculadora): \ $\sqrt{10}$, \ $\frac{19}{6}$, \ $\pi$, \ $3,\widehat1$, \ $3,14$.


% =====================================================
\seccio{2}{Decimales y fracción generatriz}

\exf{1} Indica si cada decimal es exacto, periódico puro, periódico mixto o no periódico, y si es racional o irracional:\begin{tasks}(3)
  \task $3,25$
  \task $1,\widehat{36}$
  \task $0,8\widehat3$
  \task $2,7182818\ldots$ (el número $e$)
  \task $5,\widehat2$
  \task $0,12\widehat{45}$
\end{tasks}

\ex{1} Calcula la fracción generatriz y simplifícala:\begin{tasks}(3)
  \task $0,75$
  \task $1,\widehat4$
  \task $2,\widehat{18}$
  \task $0,1\widehat6$
  \task $3,2\widehat5$
  \task $0,0\widehat3$
\end{tasks}\recorda{Periódico puro: $2,\widehat7=\frac{27-2}{9}$. \ Periódico mixto: $1,2\widehat3=\frac{123-12}{90}$ (un 9 por cada cifra del periodo y un 0 por cada cifra del anteperiodo).}

\ex{2} Escribe los decimales en forma de fracción y opera:\begin{tasks}(2)
  \task $0,\widehat3+0,\widehat6$
  \task $2,\widehat5-1,\widehat2$
  \task $1,\widehat2\cdot0,\widehat6$
  \task $0,1\widehat6\cdot1,2$
\end{tasks}

\ex{3} Calcula la fracción generatriz de $0,\widehat9$ y de $2,4\widehat9$. ¿Qué observas?


% =====================================================
\seccio{3}{Intervalos y semirrectas}

\exf{1} Escribe como intervalo o semirrecta y represéntalo en la recta real:\begin{tasks}(3)
  \task $-2\le x\le3$
  \task $1<x<5$
  \task $x\ge-1$
  \task $x<4$
  \task $-3<x\le0$
  \task $x\le\frac52$
\end{tasks}\recorda{Si el extremo pertenece ($\le$, $\ge$), corchete $[\ ]$ y punto lleno; si no ($<$, $>$), paréntesis $(\ )$ y punto vacío. En el infinito, siempre paréntesis.}

\ex{1} Escribe la desigualdad que corresponde a cada intervalo o semirrecta y represéntalo:\begin{tasks}(4)
  \task $[-4,\, 1)$
  \task $(0,\, 6]$
  \task $(-\infty,\, -2]$
  \task $(3,\, +\infty)$
\end{tasks}

\ex{1} Escribe el intervalo o la semirrecta que hay representado en cada caso, y también la desigualdad:\begin{tasks}(2)
  \task \begin{tikzpicture}[x=0.420cm,baseline=-3pt]
\draw[-latex,muted] (-5.6,0)--(6.7,0);
\draw[muted] (-5,0.1)--(-5,-0.1) node[below=-1pt,font=\scriptsize,text=muted]{$-5$};
\draw[muted] (-4,0.1)--(-4,-0.1) node[below=-1pt,font=\scriptsize,text=muted]{$-4$};
\draw[muted] (-3,0.1)--(-3,-0.1) node[below=-1pt,font=\scriptsize,text=muted]{$-3$};
\draw[muted] (-2,0.1)--(-2,-0.1) node[below=-1pt,font=\scriptsize,text=muted]{$-2$};
\draw[muted] (-1,0.1)--(-1,-0.1) node[below=-1pt,font=\scriptsize,text=muted]{$-1$};
\draw[muted] (0,0.1)--(0,-0.1) node[below=-1pt,font=\scriptsize,text=muted]{$0$};
\draw[muted] (1,0.1)--(1,-0.1) node[below=-1pt,font=\scriptsize,text=muted]{$1$};
\draw[muted] (2,0.1)--(2,-0.1) node[below=-1pt,font=\scriptsize,text=muted]{$2$};
\draw[muted] (3,0.1)--(3,-0.1) node[below=-1pt,font=\scriptsize,text=muted]{$3$};
\draw[muted] (4,0.1)--(4,-0.1) node[below=-1pt,font=\scriptsize,text=muted]{$4$};
\draw[muted] (5,0.1)--(5,-0.1) node[below=-1pt,font=\scriptsize,text=muted]{$5$};
\draw[muted] (6,0.1)--(6,-0.1) node[below=-1pt,font=\scriptsize,text=muted]{$6$};
\draw[accent,line width=2.6pt] (-3,0)--(1,0);
\fill[accent] (-3,0) circle (2.3pt);
\filldraw[fill=white,draw=accent,line width=1.1pt] (1,0) circle (2.3pt);
\end{tikzpicture}
  \task \begin{tikzpicture}[x=0.420cm,baseline=-3pt]
\draw[-latex,muted] (-5.6,0)--(6.7,0);
\draw[muted] (-5,0.1)--(-5,-0.1) node[below=-1pt,font=\scriptsize,text=muted]{$-5$};
\draw[muted] (-4,0.1)--(-4,-0.1) node[below=-1pt,font=\scriptsize,text=muted]{$-4$};
\draw[muted] (-3,0.1)--(-3,-0.1) node[below=-1pt,font=\scriptsize,text=muted]{$-3$};
\draw[muted] (-2,0.1)--(-2,-0.1) node[below=-1pt,font=\scriptsize,text=muted]{$-2$};
\draw[muted] (-1,0.1)--(-1,-0.1) node[below=-1pt,font=\scriptsize,text=muted]{$-1$};
\draw[muted] (0,0.1)--(0,-0.1) node[below=-1pt,font=\scriptsize,text=muted]{$0$};
\draw[muted] (1,0.1)--(1,-0.1) node[below=-1pt,font=\scriptsize,text=muted]{$1$};
\draw[muted] (2,0.1)--(2,-0.1) node[below=-1pt,font=\scriptsize,text=muted]{$2$};
\draw[muted] (3,0.1)--(3,-0.1) node[below=-1pt,font=\scriptsize,text=muted]{$3$};
\draw[muted] (4,0.1)--(4,-0.1) node[below=-1pt,font=\scriptsize,text=muted]{$4$};
\draw[muted] (5,0.1)--(5,-0.1) node[below=-1pt,font=\scriptsize,text=muted]{$5$};
\draw[muted] (6,0.1)--(6,-0.1) node[below=-1pt,font=\scriptsize,text=muted]{$6$};
\draw[accent,line width=2.6pt] (-5.45,0)--(2,0);
\fill[accent] (-5.75,0)--(-5.35,0.17)--(-5.35,-0.17)--cycle;
\fill[accent] (2,0) circle (2.3pt);
\end{tikzpicture}
  \task \begin{tikzpicture}[x=0.420cm,baseline=-3pt]
\draw[-latex,muted] (-5.6,0)--(6.7,0);
\draw[muted] (-5,0.1)--(-5,-0.1) node[below=-1pt,font=\scriptsize,text=muted]{$-5$};
\draw[muted] (-4,0.1)--(-4,-0.1) node[below=-1pt,font=\scriptsize,text=muted]{$-4$};
\draw[muted] (-3,0.1)--(-3,-0.1) node[below=-1pt,font=\scriptsize,text=muted]{$-3$};
\draw[muted] (-2,0.1)--(-2,-0.1) node[below=-1pt,font=\scriptsize,text=muted]{$-2$};
\draw[muted] (-1,0.1)--(-1,-0.1) node[below=-1pt,font=\scriptsize,text=muted]{$-1$};
\draw[muted] (0,0.1)--(0,-0.1) node[below=-1pt,font=\scriptsize,text=muted]{$0$};
\draw[muted] (1,0.1)--(1,-0.1) node[below=-1pt,font=\scriptsize,text=muted]{$1$};
\draw[muted] (2,0.1)--(2,-0.1) node[below=-1pt,font=\scriptsize,text=muted]{$2$};
\draw[muted] (3,0.1)--(3,-0.1) node[below=-1pt,font=\scriptsize,text=muted]{$3$};
\draw[muted] (4,0.1)--(4,-0.1) node[below=-1pt,font=\scriptsize,text=muted]{$4$};
\draw[muted] (5,0.1)--(5,-0.1) node[below=-1pt,font=\scriptsize,text=muted]{$5$};
\draw[muted] (6,0.1)--(6,-0.1) node[below=-1pt,font=\scriptsize,text=muted]{$6$};
\draw[accent,line width=2.6pt] (-1,0)--(4,0);
\filldraw[fill=white,draw=accent,line width=1.1pt] (-1,0) circle (2.3pt);
\fill[accent] (4,0) circle (2.3pt);
\end{tikzpicture}
  \task \begin{tikzpicture}[x=0.420cm,baseline=-3pt]
\draw[-latex,muted] (-5.6,0)--(6.7,0);
\draw[muted] (-5,0.1)--(-5,-0.1) node[below=-1pt,font=\scriptsize,text=muted]{$-5$};
\draw[muted] (-4,0.1)--(-4,-0.1) node[below=-1pt,font=\scriptsize,text=muted]{$-4$};
\draw[muted] (-3,0.1)--(-3,-0.1) node[below=-1pt,font=\scriptsize,text=muted]{$-3$};
\draw[muted] (-2,0.1)--(-2,-0.1) node[below=-1pt,font=\scriptsize,text=muted]{$-2$};
\draw[muted] (-1,0.1)--(-1,-0.1) node[below=-1pt,font=\scriptsize,text=muted]{$-1$};
\draw[muted] (0,0.1)--(0,-0.1) node[below=-1pt,font=\scriptsize,text=muted]{$0$};
\draw[muted] (1,0.1)--(1,-0.1) node[below=-1pt,font=\scriptsize,text=muted]{$1$};
\draw[muted] (2,0.1)--(2,-0.1) node[below=-1pt,font=\scriptsize,text=muted]{$2$};
\draw[muted] (3,0.1)--(3,-0.1) node[below=-1pt,font=\scriptsize,text=muted]{$3$};
\draw[muted] (4,0.1)--(4,-0.1) node[below=-1pt,font=\scriptsize,text=muted]{$4$};
\draw[muted] (5,0.1)--(5,-0.1) node[below=-1pt,font=\scriptsize,text=muted]{$5$};
\draw[muted] (6,0.1)--(6,-0.1) node[below=-1pt,font=\scriptsize,text=muted]{$6$};
\draw[accent,line width=2.6pt] (0,0)--(6.45,0);
\fill[accent] (6.75,0)--(6.35,0.17)--(6.35,-0.17)--cycle;
\filldraw[fill=white,draw=accent,line width=1.1pt] (0,0) circle (2.3pt);
\end{tikzpicture}
\end{tasks}

\ex{2} Expresa con un intervalo o una semirrecta:\begin{tasks}(1)
  \task los números mayores que $-3$ y menores o iguales que 2;
  \task los números reales no negativos;
  \task los números menores que 5;
  \task los números positivos menores que 1;
  \task los números que están a una distancia del 0 menor que 3.
\end{tasks}

\ex{2} ¿Qué números enteros pertenecen a cada intervalo?\begin{tasks}(4)
  \task $(-3,\, 2]$
  \task $[0,\, 4)$
  \task $(-1,\, 1)$
  \task $\left(-\frac52,\, \frac72\right)$
\end{tasks}

\ex{2} Expresa cada situación con un intervalo o una semirrecta:\begin{tasks}(1)
  \task Un ascensor admite una carga $P$ de como máximo 450 kg.
  \task Para sacarse el carné de conducir hay que tener una edad $e$ de al menos 18 años.
  \task Un termostato mantiene la temperatura $t$ por encima de 19\,$^\circ$C pero por debajo de 23\,$^\circ$C.
  \task En una prueba, aprueban las notas $n$ iguales o superiores a 5 (la nota máxima es 10).
\end{tasks}


% =====================================================
\seccio{4}{Operaciones con intervalos}

\exf{1} Dados $A=[-2,\, 3]$ y $B=(1,\, 5)$, represéntalos en la misma recta y calcula $A\cup B$ y $A\cap B$.\recorda{La unión $A\cup B$ son los números que están en $A$ \textbf{o} en $B$; la intersección $A\cap B$, los que están en $A$ \textbf{y} en $B$ a la vez.}

\ex{2} Sean $A=(-\infty,\, 1]$, $B=[-3,\, 4)$ y $C=(2,\, +\infty)$. Calcula:\begin{tasks}(3)
  \task $A\cup B$
  \task $A\cap B$
  \task $B\cap C$
  \task $A\cap C$
  \task $B\cup C$
  \task $A\cup C$
\end{tasks}

\ex{2} Calcula y representa:\begin{tasks}(3)
  \task $[0,\, 4)\cup[4,\, 7]$
  \task $[0,\, 4)\cap[4,\, 7]$
  \task $(-5,\, -1]\cap[-1,\, 2)$
  \task $(-\infty,\, 3)\cap(-\infty,\, 0]$
  \task $(-2,\, 6)\cap[1,\, 3]$
\end{tasks}

\ex{2} Escribe con un intervalo o una semirrecta los números que cumplen a la vez las dos condiciones:\begin{tasks}(2)
  \task $x>-2$ \ y \ $x\le5$
  \task $x\ge1$ \ y \ $x<1$
  \task $x\le4$ \ y \ $x\le-1$
  \task $x>0$ \ y \ $x\ge3$
\end{tasks}

\ex{2} \begin{tasks}(1)
  \task Un medicamento se tiene que conservar entre 2\,$^\circ$C y 8\,$^\circ$C (ambos incluidos). La cámara de un camión mantiene la temperatura entre 5\,$^\circ$C y 12\,$^\circ$C. Exprésalo con intervalos. ¿Entre qué temperaturas puede estar el medicamento dentro del camión?
  \task Un museo hace entrada reducida a los menores de 12 años y a los mayores de 65. Expresa con intervalos las edades que tienen entrada reducida.
\end{tasks}

\ex{3} Halla dos intervalos $A$ y $B$ tales que $A\cap B=[2,\, 3]$ y $A\cup B=(-1,\, 6)$.


% =====================================================
\seccio{5}{Raíces y radicales}

\exf{1} Calcula las raíces siguientes, si existen:\begin{tasks}(4)
  \task $\sqrt{81}$
  \task $\sqrt[3]{64}$
  \task $\sqrt[4]{16}$
  \task $\sqrt[5]{-1}$
  \task $\sqrt[3]{-1000}$
  \task $\sqrt{0,49}$
  \task $\sqrt[3]{\frac{27}{8}}$
  \task $\sqrt[4]{-81}$
\end{tasks}\recorda{$\sqrt[n]{a}=x\iff x^n=a$. Si el índice es par y el radicando negativo, la raíz no es un número real.}

\ex{1} ¿Cuántas raíces reales tiene cada radical? Calcúlalas:\begin{tasks}(5)
  \task $\sqrt{36}$
  \task $\sqrt[3]{-125}$
  \task $\sqrt[4]{-1}$
  \task $\sqrt[6]{64}$
  \task $\sqrt[5]{32}$
\end{tasks}

\ex{2} Sin calculadora, di entre qué dos números enteros consecutivos está cada raíz:\begin{tasks}(5)
  \task $\sqrt{30}$
  \task $\sqrt[3]{50}$
  \task $\sqrt{90}$
  \task $\sqrt[4]{100}$
  \task $\sqrt[3]{-20}$
\end{tasks}

\ex{1} Con la calculadora, aproxima a las centésimas:\begin{tasks}(4)
  \task $\sqrt7$
  \task $\sqrt[3]{20}$
  \task $\sqrt[4]{50}$
  \task $\sqrt[5]{-100}$
\end{tasks}


% =====================================================
\seccio{6}{Radicales y potencias de exponente fraccionario}

\exf{1} Escribe en forma de potencia:\begin{tasks}(6)
  \task $\sqrt3$
  \task $\sqrt[3]{5^2}$
  \task $\sqrt[4]{x^3}$
  \task $\sqrt[5]2$
  \task $\sqrt{a^5}$
  \task $\sqrt[3]7$
\end{tasks}\recorda{$\sqrt[n]{a^m}=a^{\frac mn}$: el exponente del radicando va al numerador y el índice, al denominador.}

\ex{1} Escribe en forma de radical:\begin{tasks}(6)
  \task $5^{\frac12}$
  \task $2^{\frac34}$
  \task $3^{\frac25}$
  \task $a^{\frac53}$
  \task $10^{\frac13}$
  \task $x^{\frac72}$
\end{tasks}

\ex{2} Escribe como una potencia de un número primo:\begin{tasks}(6)
  \task $\sqrt8$
  \task $\sqrt[3]9$
  \task $\dfrac{1}{\sqrt5}$
  \task $\dfrac{1}{\sqrt[3]4}$
  \task $\sqrt[4]{27}$
  \task $\sqrt{\dfrac12}$
\end{tasks}

\ex{2} Calcula sin calculadora:\begin{tasks}(4)
  \task $9^{\frac12}$
  \task $8^{\frac13}$
  \task $16^{\frac34}$
  \task $27^{\frac23}$
  \task $32^{\frac25}$
  \task $4^{-\frac12}$
  \task $125^{-\frac23}$
  \task $100^{\frac32}$
\end{tasks}\recorda{Escribe la base como potencia: $16^{\frac34}=\left(2^4\right)^{\frac34}=2^3$.}

\ex{3} Expresa como una sola potencia y después como un radical:\begin{tasks}(3)
  \task $2^{\frac12}\cdot2^{\frac13}$
  \task $3^{\frac34}:3^{\frac12}$
  \task $\left(5^{\frac23}\right)^{\frac34}$
  \task $\sqrt3\cdot\sqrt[3]3$
  \task $\sqrt[3]{a^2}\cdot\sqrt a$
  \task $\dfrac{\sqrt x}{\sqrt[4]x}$
\end{tasks}\recorda{$a^p\cdot a^q=a^{p+q}$, \ $a^p:a^q=a^{p-q}$, \ $\left(a^p\right)^q=a^{p\cdot q}$.}


% =====================================================
\seccio{7}{Radicales equivalentes}

\exf{1} Escribe tres radicales equivalentes a cada uno de estos:\begin{tasks}(3)
  \task $\sqrt2$
  \task $\sqrt[3]5$
  \task $\sqrt[5]{a^2}$
\end{tasks}\recorda{$\sqrt[n]{a^m}=\sqrt[n\cdot p]{a^{m\cdot p}}$: multiplicamos el índice y el exponente por el mismo número.}

\ex{1} Di si son equivalentes los radicales de cada pareja. Justifícalo pasándolos a potencia:\begin{tasks}(2)
  \task $\sqrt3$ \ y \ $\sqrt[4]9$
  \task $\sqrt[3]2$ \ y \ $\sqrt[6]8$
  \task $\sqrt[6]{25}$ \ y \ $\sqrt[3]5$
  \task $\sqrt[4]{a^3}$ \ y \ $\sqrt[8]{a^6}$
\end{tasks}

\ex{2} Simplifica los radicales siguientes:\begin{tasks}(4)
  \task $\sqrt[4]9$
  \task $\sqrt[6]8$
  \task $\sqrt[8]{16}$
  \task $\sqrt[9]{27}$
  \task $\sqrt[10]{32}$
  \task $\sqrt[6]{a^4}$
  \task $\sqrt[12]{x^9}$
  \task $\sqrt[6]{a^2b^4}$
\end{tasks}\recorda{Descompón el radicando y divide el índice y los exponentes por su m.c.d.}

\ex{2} Reduce a índice común:\begin{tasks}(3)
  \task $\sqrt2$ \ y \ $\sqrt[3]3$
  \task $\sqrt5$ \ y \ $\sqrt[4]3$
  \task $\sqrt[3]2$, \ $\sqrt[4]5$ \ y \ $\sqrt[6]3$
\end{tasks}

\ex{3} Ordena de menor a mayor, reduciendo primero a índice común:\begin{tasks}(2)
  \task $\sqrt2$, \ $\sqrt[3]3$, \ $\sqrt[6]6$
  \task $\sqrt[3]4$, \ $\sqrt3$, \ $\sqrt[4]5$
\end{tasks}


% =====================================================
\seccio{8}{Extracción e introducción de factores}

\exf{1} Extrae todos los factores posibles:\begin{tasks}(5)
  \task $\sqrt{12}$
  \task $\sqrt{18}$
  \task $\sqrt{50}$
  \task $\sqrt{75}$
  \task $\sqrt{98}$
  \task $\sqrt{200}$
  \task $\sqrt[3]{16}$
  \task $\sqrt[3]{54}$
  \task $\sqrt[3]{81}$
  \task $\sqrt[3]{250}$
\end{tasks}\recorda{Descompón el radicando en factores primos. Sale fuera cada factor que tenga exponente igual o mayor que el índice: $\sqrt{72}=\sqrt{2^3\cdot3^2}=2\cdot3\sqrt2=6\sqrt2$.}

\ex{2} Extrae todos los factores posibles (las letras representan números positivos):\begin{tasks}(3)
  \task $\sqrt{a^3}$
  \task $\sqrt{8x^5}$
  \task $\sqrt[3]{a^4b^6}$
  \task $\sqrt{72a^2b^3}$
  \task $\sqrt[3]{24x^7}$
  \task $\sqrt[4]{32a^9}$
\end{tasks}

\ex{1} Introduce los factores dentro del radical:\begin{tasks}(4)
  \task $2\sqrt3$
  \task $3\sqrt2$
  \task $5\sqrt5$
  \task $2\sqrt[3]5$
  \task $3\sqrt[3]2$
  \task $a\sqrt a$
  \task $2x\sqrt x$
\end{tasks}\recorda{El factor entra elevado al índice: $a\sqrt[n]{b}=\sqrt[n]{a^n\cdot b}$.}

\ex{2} Introduce los factores y di qué número es mayor en cada pareja:\begin{tasks}(2)
  \task $3\sqrt5$ \ o \ $5\sqrt3$
  \task $2\sqrt[3]3$ \ o \ $3\sqrt[3]2$
\end{tasks}

\ex{3} Extrae factores (haz la raíz del numerador y la del denominador):\begin{tasks}(4)
  \task $\sqrt{\dfrac{18}{25}}$
  \task $\sqrt{\dfrac{12}{49}}$
  \task $\sqrt[3]{\dfrac{16}{27}}$
  \task $\sqrt{\dfrac{75}{4}}$
\end{tasks}


% =====================================================
\seccio{9}{Operaciones con radicales}

\exf{1} Opera los radicales semejantes:\begin{tasks}(2)
  \task $3\sqrt2+5\sqrt2-\sqrt2$
  \task $4\sqrt[3]5-7\sqrt[3]5$
  \task $2\sqrt3+\sqrt3-6\sqrt3$
  \task $\sqrt7+3\sqrt5-4\sqrt7+\sqrt5$
\end{tasks}\recorda{Solo se pueden sumar o restar radicales semejantes (mismo índice y mismo radicando): se suman los coeficientes. $\sqrt2+\sqrt3$ no se puede juntar.}

\ex{2} Extrae factores y opera:\begin{tasks}(2)
  \task $\sqrt8+\sqrt{18}$
  \task $\sqrt{12}+\sqrt{27}-\sqrt3$
  \task $\sqrt{50}-\sqrt{32}+\sqrt2$
  \task $2\sqrt{45}-\sqrt{20}+\sqrt5$
  \task $3\sqrt{12}-2\sqrt{27}+\sqrt{75}$
  \task $\sqrt[3]{16}+\sqrt[3]{54}-\sqrt[3]2$
  \task $\sqrt{20}+\sqrt{45}-\sqrt{80}$
  \task $5\sqrt8-2\sqrt{18}+\sqrt{72}$
\end{tasks}

\ex{1} Multiplica o divide y simplifica el resultado:\begin{tasks}(5)
  \task $\sqrt2\cdot\sqrt8$
  \task $\sqrt3\cdot\sqrt{12}$
  \task $\sqrt5\cdot\sqrt{10}$
  \task $2\sqrt3\cdot4\sqrt6$
  \task $\sqrt[3]4\cdot\sqrt[3]2$
  \task $\sqrt[3]9\cdot\sqrt[3]3$
  \task $\dfrac{\sqrt{50}}{\sqrt2}$
  \task $\dfrac{\sqrt{72}}{\sqrt8}$
  \task $\dfrac{\sqrt[3]{48}}{\sqrt[3]6}$
  \task $\dfrac{6\sqrt{15}}{3\sqrt5}$
\end{tasks}\recorda{$\sqrt[n]a\cdot\sqrt[n]b=\sqrt[n]{a\cdot b}$ \quad y \quad $\sqrt[n]a:\sqrt[n]b=\sqrt[n]{a:b}$ \quad (mismo índice).}

\ex{2} Calcula las potencias y escribe como un solo radical las raíces de raíces:\begin{tasks}(4)
  \task $\left(\sqrt3\right)^2$
  \task $\left(\sqrt5\right)^4$
  \task $\left(2\sqrt3\right)^2$
  \task $\left(\sqrt[3]2\right)^6$
  \task $\left(\sqrt[3]4\right)^2$
  \task $\sqrt{\sqrt3}$
  \task $\sqrt[3]{\sqrt5}$
  \task $\sqrt{\sqrt[3]{\sqrt2}}$
\end{tasks}\recorda{$\left(\sqrt[n]a\right)^p=\sqrt[n]{a^p}$ \quad y \quad $\sqrt[m]{\sqrt[n]a}=\sqrt[m\cdot n]a$.}

\ex{2} Escribe como un solo radical (reduce primero a índice común o pasa a potencia):\begin{tasks}(5)
  \task $\sqrt2\cdot\sqrt[3]2$
  \task $\sqrt3\cdot\sqrt[3]2$
  \task $\sqrt[3]5\cdot\sqrt5$
  \task $\dfrac{\sqrt2}{\sqrt[3]2}$
  \task $\dfrac{\sqrt[3]9}{\sqrt3}$
\end{tasks}

\ex{2} Aplica la propiedad distributiva y simplifica:\begin{tasks}(2)
  \task $\sqrt2\left(\sqrt2+\sqrt8\right)$
  \task $\sqrt3\left(2\sqrt3-\sqrt{12}\right)$
  \task $\left(\sqrt5+1\right)\left(\sqrt5+2\right)$
  \task $\left(2\sqrt3+1\right)\left(\sqrt3-2\right)$
\end{tasks}

\ex{3} Desarrolla con las identidades notables:\begin{tasks}(2)
  \task $\left(\sqrt3+1\right)^2$
  \task $\left(\sqrt5-\sqrt2\right)^2$
  \task $\left(\sqrt7+2\right)\left(\sqrt7-2\right)$
  \task $\left(3\sqrt2-1\right)^2$
  \task $\left(2\sqrt5+3\right)\left(2\sqrt5-3\right)$
  \task $\left(1+\sqrt2\right)^2-\left(1-\sqrt2\right)^2$
\end{tasks}\recorda{$(a+b)^2=a^2+2ab+b^2$, \ $(a-b)^2=a^2-2ab+b^2$, \ $(a+b)(a-b)=a^2-b^2$.}


% =====================================================
\seccio{10}{Racionalización}

\exf{1} Racionaliza y simplifica:\begin{tasks}(6)
  \task $\dfrac{1}{\sqrt2}$
  \task $\dfrac{3}{\sqrt3}$
  \task $\dfrac{6}{\sqrt2}$
  \task $\dfrac{5}{\sqrt{10}}$
  \task $\dfrac{2}{3\sqrt5}$
  \task $\dfrac{\sqrt2}{\sqrt6}$
\end{tasks}\recorda{Si en el denominador hay $\sqrt b$, multiplica el numerador y el denominador por $\sqrt b$: $\ \dfrac{6}{\sqrt3}=\dfrac{6\sqrt3}{3}=2\sqrt3$.}

\ex{2} Racionaliza y simplifica:\begin{tasks}(6)
  \task $\dfrac{1}{\sqrt[3]2}$
  \task $\dfrac{6}{\sqrt[3]3}$
  \task $\dfrac{4}{\sqrt[3]4}$
  \task $\dfrac{10}{\sqrt[3]{25}}$
  \task $\dfrac{1}{\sqrt[4]2}$
  \task $\dfrac{3}{\sqrt[4]{27}}$
\end{tasks}\recorda{Con $\sqrt[n]{b^m}$ en el denominador, multiplica por $\sqrt[n]{b^{n-m}}$, porque $\sqrt[n]{b^m}\cdot\sqrt[n]{b^{n-m}}=\sqrt[n]{b^n}=b$.}

\ex{2} Racionaliza multiplicando por el conjugado:\begin{tasks}(4)
  \task $\dfrac{1}{\sqrt2+1}$
  \task $\dfrac{2}{\sqrt3-1}$
  \task $\dfrac{3}{2+\sqrt3}$
  \task $\dfrac{4}{\sqrt5+\sqrt3}$
  \task $\dfrac{\sqrt2}{\sqrt2-1}$
  \task $\dfrac{6}{3-\sqrt3}$
  \task $\dfrac{\sqrt3+1}{\sqrt3-1}$
  \task $\dfrac{1}{2\sqrt2-\sqrt3}$
\end{tasks}\recorda{El conjugado de $a+\sqrt b$ es $a-\sqrt b$, y $\left(a+\sqrt b\right)\left(a-\sqrt b\right)=a^2-b$.}

\ex{3} Racionaliza y opera:\begin{tasks}(2)
  \task $\dfrac{1}{\sqrt2}+\dfrac{1}{\sqrt8}$
  \task $\dfrac{2}{\sqrt3}-\sqrt3$
  \task $\dfrac{1}{\sqrt2-1}-\dfrac{1}{\sqrt2+1}$
  \task $\dfrac{\sqrt{12}}{\sqrt3+1}$
\end{tasks}


% =====================================================
\seccio{11}{Problemas}

\exf{2} Un cuadrado tiene un área de $50\ \text{cm}^2$. Calcula su lado, su perímetro y su diagonal. Da los resultados en forma de radical simplificado y, el perímetro, también aproximado a las centésimas.

\ex{2} Un rectángulo mide $\sqrt{12}$ cm de base y $\sqrt{27}$ cm de altura. Calcula su perímetro, su área y su diagonal.

\ex{2} Los catetos de un triángulo rectángulo miden $\sqrt8$ cm y $\sqrt{12}$ cm. Calcula su hipotenusa y su área.

\ex{2} Un cubo tiene un volumen de $54\ \text{cm}^3$. Calcula su arista y su superficie total en forma de radical, y aproxímalas a las centésimas.

\ex{3} Un triángulo equilátero mide $2\sqrt3$ cm de lado. Calcula su altura y su área.

\ex{3} En un folio DIN A4, el lado largo entre el lado corto es $\sqrt2$. Si el lado corto mide 21 cm, ¿cuánto mide el largo (exacto y aproximado)? Si doblamos el folio por la mitad, el nuevo lado corto es la mitad del largo: comprueba, racionalizando, que el nuevo folio también cumple que largo\,:\,corto $=\sqrt2$.

\vfill
\begin{center}\color{muted}\small\faCheckCircle[regular]\ \ ¡Ánimo! Comprueba tus resultados con las herramientas interactivas de los apuntes de la unidad.\end{center}

\end{document}
